\documentclass[12pt]{article}
\usepackage[utf8]{inputenc}
\usepackage[T1]{fontenc}
\usepackage{hyperref}
\title{Fanti Language Resources for K-12 Curriculum Development}
\author{Compiled by Claude Sonnet 5.5 (claude-sonnet-5-5) for Andr\'as Kornai}
\date{}
\begin{document}
\maketitle

\section{Introduction}
Fanti (Fante, Mfantse; ISO 639-3 \texttt{fat}) is an Akan dialect of central and southern Ghana, with a standardized written form said to be taught in primary and secondary schools. Speaker figures in the sources conflict (about 2.8 million to over 6 million). The single search run found mostly general reference pages; the list below is short.

\section{Primary Educational Resources}
No Fante primer or school textbook was found.

\section{Linguistic and Reference Materials}
\begin{itemize}
\item Wikipedia article on the Fante dialect (status, schooling, orthography committee): \url{https://en.wikipedia.org/wiki/Fante_dialect}
\item Omniglot page on Fante (alphabet chart and sample text): \url{https://omniglot.com/writing/fante.htm}
\item SOAS digital archive, Akan materials, including \emph{Mfantsi grammar} revised by J. P. Brown (2nd ed., 1913, 72 pp.): \url{https://digital.soas.ac.uk/akan/all/thumbs}
\end{itemize}

\section{Religious and Community Resources}
\begin{itemize}
\item Joshua Project Akan language page (Bible availability listed for Akan generally, not Fante specifically): \url{https://www.joshuaproject.net/index.php/languages/aka}
\end{itemize}

\section{Digital and Audio Resources}
\begin{itemize}
\item Localazy language hub entry for Fanti (software localization): \url{https://hub.localazy.com/en/languages/fat-fanti}
\end{itemize}

\section{Curriculum Development Context}
No modern Fante dictionary or classroom materials surfaced. Further searching (Ghana Education Service, Bureau of Ghana Languages, Fante Bible editions) is needed before a curriculum can be planned.
\end{document}